\begin{center}
Problema G

Zé do Grafo no Halloween

{\small Nome do arquivo: G.c, G.cpp, G.java, ou G.py}

{\large Timelimit: $1$}

{\tiny Por André da Cunha Ribeiro}

\end{center}			

Zé do Grafo, o famoso bruxo das noites de Halloween e do Abóbora Contest, adora brincar com suas poções mágicas e encantamentos geométricos. Este ano, ele está experimentando sua força com palitinhos mágicos, que têm o poder de se transformar em lados de um triângulo. No entanto, para garantir que o feitiço funcione, ele precisa seguir a regra mágica da {\bf Desigualdade Triangular}, ou seja, a soma dos comprimentos de quaisquer dois palitinhos deve ser maior que o comprimento do terceiro.

Zé do Grafo possui vários palitinhos retos de diversos tamanhos e o objetivo e formar um triângulo com eles pela rega mágica, de modo que cada palitinho represente exatamente um lado do triângulo. Não é permitido que algum lado possua um buraco, nem que um pedaço de algum palitinho sobre para fora do triângulo. A figura \ref{fig:G} ilustra um trio permitido na sub-figura (a) e as sub-figuras (b) e (c) ilustram trios proibidos de acordo com as regras mágica.

\begin{figure}[!ht]
\centering
\includegraphics[width=14cm]{fig/figuraG.png}
\caption{Exemplos de palitinhos}
\label{fig:G}
\end{figure}

Zé do Grafo quer saber se com três palitinhos distintos, ele consegue um triângulo pela regra mágica. 

\vspace{0.3cm}
\textbf{Entrada}
A entrada são três números inteiros, que representa o tamanho de cada palitinho.

\vspace{0.3cm}
\textbf{Saída}
O programa deve retornar ``SIM'' caso Zé do Grafos consiga forma o triângulo. Caso contrário, a resposta será ``NAO''. Não esqueça do terminador de linha.

\begin{table}[h]
    \centering
    \begin{tabular}{|p{7cm}|p{7cm}|}
        \hline
        \begin{center}
            \vspace{-0.3cm}
            \textbf{Exemplos de Entrada}
            \vspace{-0.3cm}
        \end{center} 
         &  
        \begin{center}
            \vspace{-0.3cm}
            \textbf{Exemplos de Saída}
            \vspace{-0.3cm}
        \end{center} \\ \hline
        \texttt{3 4 6} & \texttt{SIM} \\ \hline
        \texttt{3 4 7} & \texttt{NAO} \\ \hline
        \texttt{3 3 7} & \texttt{NAO} \\ \hline
    \end{tabular}
    \caption{Exemplos de entradas e saídas}
    
\end{table}

\newpage

\begin{center}
\noindent
\lstinputlisting{abobora_24/G/G4.cpp}
\end{center}